\documentclass[a4paper,12pt]{ctexrep}

\usepackage[top=2.5cm, bottom=2.5cm, left=2.5cm, right=2.5cm]{geometry}
\usepackage{amsmath,amssymb,amsfonts}
\usepackage{graphicx}
\usepackage{booktabs}
\usepackage{float}
\usepackage{caption}
\usepackage{enumitem}
\usepackage{multirow}
\usepackage{array}
\usepackage{tabularx}
\usepackage{subcaption}
\usepackage{bm}
\usepackage{mathrsfs}

\captionsetup[figure]{font=small, labelsep=quad, skip=6pt, belowskip=12pt}
\captionsetup[table]{font=small, labelsep=quad, skip=6pt, belowskip=6pt}

\ctexset{
    chapter = {
        format = \heiti\fontsize{16pt}{16pt}\selectfont\bfseries\centering,
        nameformat = {},
        titleformat = {},
        beforeskip = 24pt,
        afterskip = 18pt,
        lofskip = 0pt,
        lotskip = 0pt,
    },
    section = {
        format = \heiti\fontsize{14pt}{14pt}\selectfont\bfseries\centering,
        nameformat = {},
        titleformat = {},
        beforeskip = 24pt,
        afterskip = 6pt,
    },
    subsection = {
        format = \heiti\fontsize{13pt}{13pt}\selectfont\raggedright,
        nameformat = {},
        titleformat = {},
        beforeskip = 12pt,
        afterskip = 6pt,
    },
    subsubsection = {
        format = \heiti\fontsize{12pt}{12pt}\selectfont\raggedright,
        nameformat = {},
        titleformat = {},
        beforeskip = 12pt,
        afterskip = 6pt,
    }
}

\setlength{\parindent}{2em}
\setlength{\parskip}{0pt}
\linespread{1.25}

\newcommand{\song}{\songti}
\newcommand{\hei}{\heiti}
\newcommand{\kai}{\kaishu}
\newcommand{\fs}{\fangsong}

\begin{document}

% ============================================================
% 封面
% ============================================================
\begin{titlepage}
\centering
\vspace*{2cm}
{\heiti\fontsize{26pt}{26pt}\selectfont SVD与ICP算法推导\\[0.5em]及仿真实验分析}\\[3cm]
{\songti\fontsize{16pt}{16pt}\selectfont 2312167-刘振毫-智能科学与技术}\\[4cm]

{\songti\fontsize{14pt}{14pt}\selectfont 2026年5月8日}
\vfill
\end{titlepage}

% ============================================================
% 中文摘要
% ============================================================
\chapter*{摘\quad 要}
\addcontentsline{toc}{chapter}{摘\quad 要}
\markboth{摘要}{摘要}

奇异值分解（Singular Value Decomposition, SVD）是线性代数中一种重要的矩阵分解方法，它将任意实矩阵分解为两个正交矩阵和一个对角矩阵的乘积，广泛应用于数据降维、最小二乘求解和刚体变换估计等领域。迭代最近点（Iterative Closest Point, ICP）算法是三维点云配准中最为经典的方法之一，其核心思想是通过迭代寻找两组点云之间的最近点对应关系，并在每一步利用SVD求解最优刚体变换，直至收敛。本文首先从数学角度严格推导了SVD分解的存在性和构造方法，详细阐述了基于互协方差矩阵SVD分解求解最优旋转和平移的理论依据；其次，推导了ICP算法的完整迭代流程，包括最近点搜索、SVD刚体变换求解和收敛判定；然后，通过Python编程实现了SVD手动分解和ICP配准算法，并设计了四组仿真实验，分别验证了无噪声条件下的配准精度、不同噪声水平下的鲁棒性、不同旋转角度下的收敛性能以及不同重叠率下的配准效果；最后，对仿真实验结果进行了深入分析，验证了所构建模型的正确性。实验结果表明，手动实现的SVD分解与NumPy库的计算结果一致，重构误差在机器精度范围内；在无噪声条件下，ICP算法能够精确恢复真实的刚体变换参数，旋转矩阵误差和平移向量误差均在$10^{-15}$量级；配准误差随噪声水平近似线性增长，验证了算法的鲁棒性；小角度旋转收敛更快，大角度旋转可能陷入局部最优，这是ICP算法的固有局限性而非实现错误。

\vspace{1em}
\noindent\textbf{关键词：}奇异值分解；迭代最近点；点云配准；刚体变换；仿真实验

% ============================================================
% 英文摘要
% ============================================================
\chapter*{Abstract}
\addcontentsline{toc}{chapter}{Abstract}
\markboth{Abstract}{Abstract}

Singular Value Decomposition (SVD) is a fundamental matrix factorization technique in linear algebra that decomposes any real matrix into the product of two orthogonal matrices and a diagonal matrix. It is widely applied in dimensionality reduction, least squares estimation, and rigid transformation estimation. The Iterative Closest Point (ICP) algorithm is one of the most classical methods for three-dimensional point cloud registration. Its core idea is to iteratively establish the closest point correspondences between two point clouds and solve the optimal rigid transformation using SVD at each step until convergence. This report first rigorously derives the existence and construction of SVD from a mathematical perspective, elaborating the theoretical basis for solving the optimal rotation and translation via SVD of the cross-covariance matrix. Second, the complete iterative procedure of the ICP algorithm is derived, including nearest point search, SVD-based rigid transformation solving, and convergence criteria. Third, both the manual SVD decomposition and ICP registration algorithm are implemented in Python, and four groups of simulation experiments are designed to verify the registration accuracy under noise-free conditions, the robustness under different noise levels, the convergence performance under different rotation angles, and the registration effectiveness under different overlap ratios. Finally, the simulation results are thoroughly analyzed to validate the correctness of the constructed models. The experimental results demonstrate that the manually implemented SVD decomposition is consistent with the NumPy library, with reconstruction errors within machine precision. Under noise-free conditions, the ICP algorithm accurately recovers the true rigid transformation parameters, with rotation matrix errors and translation vector errors on the order of $10^{-15}$. The registration error increases approximately linearly with the noise level, confirming the algorithm's robustness. Small-angle rotations converge faster, while large-angle rotations may fall into local optima, which is an inherent limitation of the ICP algorithm rather than an implementation error.

\vspace{1em}
\noindent\textbf{Key Words:} Singular Value Decomposition; Iterative Closest Point; Point Cloud Registration; Rigid Transformation; Simulation Experiment

% ============================================================
% 目录
% ============================================================
\tableofcontents

% ============================================================
% 第一章 绪论
% ============================================================
\chapter{绪论}

三维点云配准是计算机视觉、机器人导航、医学影像和逆向工程等领域中的基础问题，其目标是寻找两组三维点云之间的空间变换关系，使得它们在公共坐标系下达到最佳对齐。在众多配准算法中，迭代最近点算法因其原理简洁、实现方便且在良好初始条件下具有较高精度而成为最为广泛使用的方法之一。ICP算法由Besl和McKay于1992年提出，其核心思想是在每次迭代中建立两组点云之间的最近点对应关系，然后基于该对应关系求解最优刚体变换，并不断重复此过程直至收敛。在ICP算法的每一步中，求解最优刚体变换的关键步骤依赖于奇异值分解，通过将互协方差矩阵进行SVD分解，可以解析地得到最优旋转矩阵和平移向量，从而避免了复杂的非线性优化。

奇异值分解作为线性代数中最重要的矩阵分解之一，其理论价值和应用价值均十分突出。从理论角度看，SVD揭示了矩阵的几何本质，即任何线性变换都可以分解为旋转、缩放和旋转三个连续操作的复合。从应用角度看，SVD在主成分分析、低秩近似、最小二乘问题和刚体变换估计等领域均有重要应用。在点云配准的背景下，SVD的作用尤为关键：它将原本需要非线性优化的刚体变换求解问题转化为可以通过矩阵分解直接求解的线性问题，极大地提高了计算效率和数值稳定性。

本文的研究目标是对SVD和ICP算法进行严格的数学推导，通过编程实现这两种算法，并设计系统的仿真实验来验证所构建模型的正确性。具体而言，本文首先从矩阵特征分解的角度推导SVD的存在性和构造方法，然后基于SVD推导ICP算法中刚体变换的最优解，接着通过Python实现这两种算法并进行四组仿真实验，最后对实验结果进行深入分析，从配准精度、噪声鲁棒性、旋转角度敏感性和重叠率依赖性等多个维度验证算法的正确性和有效性。

% ============================================================
% 第二章 SVD算法推导
% ============================================================
\chapter{SVD算法推导}

\section{SVD的定义与存在性}

奇异值分解是线性代数中一种重要的矩阵分解方法。对于任意实矩阵$\bm{A}\in\mathbb{R}^{m\times n}$，存在分解
\begin{equation}
    \bm{A} = \bm{U}\bm{\Sigma}\bm{V}^{\mathsf{T}}
\end{equation}
其中$\bm{U}\in\mathbb{R}^{m\times m}$和$\bm{V}\in\mathbb{R}^{n\times n}$均为正交矩阵，分别称为左奇异向量矩阵和右奇异向量矩阵；$\bm{\Sigma}\in\mathbb{R}^{m\times n}$为对角矩阵，其对角线元素$\sigma_1\geq\sigma_2\geq\cdots\geq\sigma_r>0$称为奇异值，$r=\mathrm{rank}(\bm{A})$为矩阵$\bm{A}$的秩。

SVD的存在性可以通过构造性证明来建立。考虑矩阵$\bm{A}^{\mathsf{T}}\bm{A}$，由于$(\bm{A}^{\mathsf{T}}\bm{A})^{\mathsf{T}}=\bm{A}^{\mathsf{T}}\bm{A}$，且对任意非零向量$\bm{x}$有$\bm{x}^{\mathsf{T}}\bm{A}^{\mathsf{T}}\bm{A}\bm{x}=\|\bm{A}\bm{x}\|^2\geq 0$，因此$\bm{A}^{\mathsf{T}}\bm{A}$是$n\times n$的对称半正定矩阵。根据实对称矩阵的谱分解定理，$\bm{A}^{\mathsf{T}}\bm{A}$可以进行正交对角化，即存在正交矩阵$\bm{V}$和对角矩阵$\bm{\Lambda}$，使得
\begin{equation}
    \bm{A}^{\mathsf{T}}\bm{A} = \bm{V}\bm{\Lambda}\bm{V}^{\mathsf{T}}
\end{equation}
其中$\bm{\Lambda}=\mathrm{diag}(\lambda_1,\lambda_2,\ldots,\lambda_n)$，$\lambda_1\geq\lambda_2\geq\cdots\geq\lambda_n\geq 0$为$\bm{A}^{\mathsf{T}}\bm{A}$的特征值。奇异值定义为特征值的非负平方根
\begin{equation}
    \sigma_i = \sqrt{\lambda_i}, \quad i=1,2,\ldots,n
\end{equation}
由于$\bm{A}^{\mathsf{T}}\bm{A}$的秩等于$\bm{A}$的秩$r$，因此恰有$r$个正特征值，对应的奇异值$\sigma_1,\sigma_2,\ldots,\sigma_r>0$，而$\sigma_{r+1}=\cdots=\sigma_n=0$。

\section{左奇异向量与右奇异向量的构造}

右奇异向量矩阵$\bm{V}$的列向量即为$\bm{A}^{\mathsf{T}}\bm{A}$的特征向量。对于左奇异向量矩阵$\bm{U}$的构造，可以利用$\bm{A}$与$\bm{V}$之间的关系。对于$\sigma_i>0$的指标$i$，定义
\begin{equation}
    \bm{u}_i = \frac{\bm{A}\bm{v}_i}{\sigma_i}, \quad i=1,2,\ldots,r
\end{equation}
可以验证这样构造的$\bm{u}_i$构成一组标准正交向量。事实上，对于$i\neq j$，有
\begin{equation}
    \bm{u}_i^{\mathsf{T}}\bm{u}_j = \frac{\bm{v}_i^{\mathsf{T}}\bm{A}^{\mathsf{T}}\bm{A}\bm{v}_j}{\sigma_i\sigma_j} = \frac{\lambda_j\bm{v}_i^{\mathsf{T}}\bm{v}_j}{\sigma_i\sigma_j} = 0
\end{equation}
对于$i=j$，有
\begin{equation}
    \bm{u}_i^{\mathsf{T}}\bm{u}_i = \frac{\bm{v}_i^{\mathsf{T}}\bm{A}^{\mathsf{T}}\bm{A}\bm{v}_i}{\sigma_i^2} = \frac{\lambda_i}{\sigma_i^2} = 1
\end{equation}
因此$\bm{u}_1,\bm{u}_2,\ldots,\bm{u}_r$确实构成一组标准正交向量。对于$\sigma_i=0$的部分，需要补充$m-r$个向量$\bm{u}_{r+1},\ldots,\bm{u}_m$，使其与$\bm{u}_1,\ldots,\bm{u}_r$共同构成$\mathbb{R}^m$的一组标准正交基。这一步骤可以通过改进的Gram-Schmidt正交化方法来实现：从随机向量出发，依次减去其在已有正交基上的投影分量，然后归一化，即可得到与已有向量正交的新基向量。

\section{宽扁矩阵的处理}

当$m<n$时，矩阵$\bm{A}$为宽扁矩阵，此时$\bm{A}^{\mathsf{T}}\bm{A}$的维度为$n\times n$，大于$\bm{A}\bm{A}^{\mathsf{T}}$的维度$m\times m$。为了减少计算量，可以转而计算$\bm{A}\bm{A}^{\mathsf{T}}$的特征分解
\begin{equation}
    \bm{A}\bm{A}^{\mathsf{T}} = \bm{U}\bm{\Lambda}'\bm{U}^{\mathsf{T}}
\end{equation}
此时左奇异向量矩阵$\bm{U}$即为$\bm{A}\bm{A}^{\mathsf{T}}$的特征向量矩阵，而右奇异向量可以通过关系式
\begin{equation}
    \bm{v}_i = \frac{\bm{A}^{\mathsf{T}}\bm{u}_i}{\sigma_i}, \quad i=1,2,\ldots,r
\end{equation}
来构造，同样需要用改进的Gram-Schmidt方法补充正交基向量。这种根据矩阵形状选择不同计算路径的策略，使得SVD分解的数值计算更加高效和稳定。

\section{SVD的几何意义}

从几何角度看，SVD揭示了线性变换$\bm{x}\mapsto\bm{A}\bm{x}$的本质结构。正交矩阵$\bm{V}^{\mathsf{T}}$表示在输入空间（$\mathbb{R}^n$）中的一个旋转或反射，对角矩阵$\bm{\Sigma}$表示沿各坐标轴的缩放（缩放因子即为奇异值），正交矩阵$\bm{U}$表示在输出空间（$\mathbb{R}^m$）中的一个旋转或反射。因此，任何线性变换都可以分解为"旋转—缩放—旋转"三个步骤的复合。这一几何解释不仅加深了对SVD的理解，也为后续在刚体变换求解中的应用奠定了理论基础。

% ============================================================
% 第三章 ICP算法推导
% ============================================================
\chapter{ICP算法推导}

\section{问题描述}

给定两组三维点云$\mathcal{P}=\{\bm{p}_i\}_{i=1}^{N}$和$\mathcal{Q}=\{\bm{q}_i\}_{i=1}^{N}$，ICP算法的目标是求解最优刚体变换（旋转矩阵$\bm{R}$和平移向量$\bm{t}$），使得目标函数
\begin{equation}
    f(\bm{R},\bm{t}) = \sum_{i=1}^{N}\|\bm{R}\bm{p}_i+\bm{t}-\bm{q}_i\|^2
\end{equation}
达到最小值，其中$\bm{R}\in\mathrm{SO}(3)$满足正交性约束$\bm{R}^{\mathsf{T}}\bm{R}=\bm{I}$和行列式约束$\det(\bm{R})=1$，$\bm{t}\in\mathbb{R}^3$。

\section{平移向量的解耦}

该优化问题的关键观察在于，平移向量$\bm{t}$可以通过质心关系从目标函数中消去。定义两组点云的质心为
\begin{equation}
    \bar{\bm{p}} = \frac{1}{N}\sum_{i=1}^{N}\bm{p}_i, \quad \bar{\bm{q}} = \frac{1}{N}\sum_{i=1}^{N}\bm{q}_i
\end{equation}
将目标函数关于$\bm{t}$求导并令其等于零，可得最优平移向量满足
\begin{equation}
    \bm{t}^* = \bar{\bm{q}} - \bm{R}\bar{\bm{p}}
\end{equation}
将上式代回目标函数，并定义去质心化的点集
\begin{equation}
    \bm{p}'_i = \bm{p}_i - \bar{\bm{p}}, \quad \bm{q}'_i = \bm{q}_i - \bar{\bm{q}}
\end{equation}
则目标函数简化为仅关于旋转矩阵$\bm{R}$的优化问题
\begin{equation}
    \min_{\bm{R}} \sum_{i=1}^{N}\|\bm{R}\bm{p}'_i - \bm{q}'_i\|^2
\end{equation}

\section{基于SVD的最优旋转求解}

展开上述目标函数，有
\begin{equation}
    \sum_{i=1}^{N}\|\bm{R}\bm{p}'_i - \bm{q}'_i\|^2 = \sum_{i=1}^{N}(\|\bm{p}'_i\|^2 + \|\bm{q}'_i\|^2) - 2\sum_{i=1}^{N}\bm{q}'^{\mathsf{T}}_i\bm{R}\bm{p}'_i
\end{equation}
由于第一项与$\bm{R}$无关，最小化目标函数等价于最大化
\begin{equation}
    \sum_{i=1}^{N}\bm{q}'^{\mathsf{T}}_i\bm{R}\bm{p}'_i = \mathrm{tr}\left(\bm{R}\sum_{i=1}^{N}\bm{p}'_i\bm{q}'^{\mathsf{T}}_i\right) = \mathrm{tr}(\bm{R}\bm{H})
\end{equation}
其中$\bm{H}=\sum_{i=1}^{N}\bm{p}'_i\bm{q}'^{\mathsf{T}}_i$为互协方差矩阵。若将去质心化后的点集表示为矩阵$\bm{P}'\in\mathbb{R}^{N\times 3}$和$\bm{Q}'\in\mathbb{R}^{N\times 3}$，则互协方差矩阵可以紧凑地写为$\bm{H}=\bm{P}'^{\mathsf{T}}\bm{Q}'$。

对$\bm{H}$进行SVD分解，设$\bm{H}=\bm{U}\bm{\Sigma}\bm{V}^{\mathsf{T}}$，则
\begin{equation}
    \mathrm{tr}(\bm{R}\bm{H}) = \mathrm{tr}(\bm{R}\bm{U}\bm{\Sigma}\bm{V}^{\mathsf{T}}) = \mathrm{tr}(\bm{V}^{\mathsf{T}}\bm{R}\bm{U}\bm{\Sigma})
\end{equation}
令$\bm{Z}=\bm{V}^{\mathsf{T}}\bm{R}\bm{U}$，由于$\bm{R}$、$\bm{U}$、$\bm{V}$均为正交矩阵，$\bm{Z}$也是正交矩阵。设$\bm{Z}$的对角线元素为$z_{ii}$，由于正交矩阵的每个元素绝对值不超过1，有
\begin{equation}
    \mathrm{tr}(\bm{Z}\bm{\Sigma}) = \sum_{i=1}^{3}z_{ii}\sigma_i \leq \sum_{i=1}^{3}\sigma_i
\end{equation}
等号成立当且仅当$z_{ii}=1$，即$\bm{Z}=\bm{I}$。因此最优旋转矩阵为
\begin{equation}
    \bm{R}^* = \bm{V}\bm{U}^{\mathsf{T}}
\end{equation}

\section{反射修正}

上述推导得到的$\bm{R}^*=\bm{V}\bm{U}^{\mathsf{T}}$满足正交性，但不一定满足$\det(\bm{R}^*)=1$。当$\det(\bm{R}^*)=-1$时，$\bm{R}^*$表示一个反射而非旋转，需要修正。修正方法是将$\bm{V}$的最后一列取反，即令$\bm{V}'=\bm{V}\cdot\mathrm{diag}(1,1,-1)$，然后计算$\bm{R}^*=\bm{V}'\bm{U}^{\mathsf{T}}$。这一修正的依据在于，当$\det(\bm{V}\bm{U}^{\mathsf{T}})=-1$时，$\bm{Z}$的行列式为$-1$，这意味着$\bm{Z}$中恰有一个对角线元素为$-1$，其余为$1$。将$\bm{V}$的最后一列取反等价于将$\bm{Z}$的第三行取反，使得$\bm{Z}$变为单位矩阵，此时$\mathrm{tr}(\bm{Z}\bm{\Sigma})=\sigma_1+\sigma_2-\sigma_3$，在$\sigma_3$为最小奇异值的情况下，这是满足$\det(\bm{R})=1$约束下的最大值。

\section{ICP迭代流程}

ICP算法的完整迭代流程如下：首先，对源点云中的每个点，在目标点云中寻找距离最近的点，建立点对对应关系；其次，基于当前的对应关系，利用上述SVD方法求解最优旋转矩阵和平移向量；然后，用求得的刚体变换更新源点云的位置；最后，计算当前均方误差，若误差变化小于预设阈值则停止迭代，否则返回第一步继续迭代。ICP算法的收敛性由以下两个事实保证：第一，在固定对应关系下，SVD给出了最优刚体变换的解析解，因此每步都使目标函数不增；第二，目标函数有下界（非负），因此序列必然收敛。但需要注意的是，ICP仅保证收敛到局部最小值，初始位姿的选择对最终配准结果有重要影响。

% ============================================================
% 第四章 仿真实验
% ============================================================
\chapter{仿真实验}

\section{实验环境与参数设置}

本实验使用Python 3.12编程语言，依赖NumPy 2.4.4进行矩阵运算和Matplotlib 3.10.9进行可视化。实验中生成三维点云数据作为测试对象，采用模拟兔子形状的非对称点云，包含椭球体身体和两个耳朵部分，共计300个点。真实刚体变换设置为绕轴$(0.3, 0.7, 0.6)$（归一化后）旋转$35°$，平移向量为$(1.5, -0.8, 0.6)$。ICP算法的最大迭代次数设为100，收敛阈值设为$10^{-10}$。

\section{实验一：无噪声条件下的ICP配准}

在无噪声条件下，源点云和目标点云之间仅存在已知的刚体变换关系，不存在任何测量误差或随机扰动。ICP算法经过26次迭代后收敛，最终均方误差从初始的1.924降至接近零。估计的旋转角度为$34.9999°$，与真实旋转角度$35°$之间的误差小于$0.001°$；旋转矩阵的Frobenius误差为$2.77\times 10^{-15}$，平移向量误差为$5.98\times 10^{-16}$，均在机器精度范围内。估计旋转矩阵的正交性误差为$5.39\times 10^{-15}$，行列式为1.0000000000，满足刚体变换的约束条件。这一结果验证了ICP算法在理想条件下能够精确恢复真实的刚体变换参数，同时也验证了基于SVD的刚体变换求解方法的正确性。

\begin{figure}[H]
    \centering
    \includegraphics[width=0.95\textwidth]{icp_registration_3d.png}
    \caption{ICP配准前后点云对比及配准后点对距离分布}
    \label{fig:registration_3d}
\end{figure}

图~\ref{fig:registration_3d}展示了ICP配准前后的三维点云对比。左图为配准前的源点云（蓝色）和目标点云（红色），两者在空间中存在明显的旋转和平移偏移；中图为配准后的结果，源点云经变换后与目标点云高度重合；右图为配准后点对距离的分布直方图，距离值集中在极小值附近，进一步证实了配准的高精度。

\section{实验二：不同噪声水平下的ICP配准}

为评估ICP算法对噪声的鲁棒性，在源点云和目标点云上分别添加不同标准差的高斯白噪声，噪声标准差$\sigma$取值为0.00、0.01、0.02、0.05、0.10和0.20。表~\ref{tab:noise}列出了不同噪声水平下的配准误差。

\begin{table}[H]
    \centering
    \caption{不同噪声水平下的ICP配准误差}
    \label{tab:noise}
    \begin{tabular}{cccc}
        \toprule
        噪声标准差$\sigma$ & 旋转误差 & 平移误差 & 迭代次数 \\
        \midrule
        0.000 & $2.44\times 10^{-15}$ & $1.11\times 10^{-15}$ & 27 \\
        0.010 & $1.54\times 10^{-3}$  & $1.28\times 10^{-3}$  & 42 \\
        0.020 & $1.32\times 10^{-3}$  & $8.05\times 10^{-4}$  & 31 \\
        0.050 & $1.74\times 10^{-2}$  & $1.12\times 10^{-2}$  & 30 \\
        0.100 & $2.50\times 10^{-2}$  & $2.45\times 10^{-2}$  & 21 \\
        0.200 & $1.06\times 10^{-1}$  & $1.08\times 10^{-1}$  & 30 \\
        \bottomrule
    \end{tabular}
\end{table}

从表~\ref{tab:noise}可以看出，当$\sigma=0$时，配准误差在机器精度范围内；随着噪声水平的增大，旋转误差和平移误差均呈增长趋势。通过线性回归分析，旋转误差随噪声的增长斜率约为0.543，平移误差的增长斜率约为0.564，两者均接近0.5，表明配准误差与噪声水平近似呈线性关系。这一结果符合理论预期：在ICP的每次迭代中，噪声导致最近点对应关系存在偏差，而SVD求解的刚体变换是当前对应关系下的最优解，因此最终的配准精度受限于对应关系的准确性，而对应关系的准确性又直接受噪声影响。值得注意的是，即使在$\sigma=0.2$的较高噪声水平下，旋转误差仍为0.106，平移误差为0.108，说明ICP算法对中等程度的噪声具有良好的鲁棒性。

\section{实验三：不同旋转角度下的ICP配准}

为研究初始旋转角度对ICP配准性能的影响，将真实旋转角度分别设置为$5°$、$15°$、$30°$、$45°$、$60°$和$90°$，平移向量保持不变。表~\ref{tab:angle}列出了不同旋转角度下的配准结果。

\begin{table}[H]
    \centering
    \caption{不同旋转角度下的ICP配准结果}
    \label{tab:angle}
    \begin{tabular}{cccc}
        \toprule
        旋转角度 & 旋转误差 & 平移误差 & 迭代次数 \\
        \midrule
        $5°$  & $2.53\times 10^{-15}$ & $1.37\times 10^{-15}$ & 17 \\
        $15°$ & $2.40\times 10^{-15}$ & $8.95\times 10^{-16}$ & 19 \\
        $30°$ & $5.44\times 10^{-15}$ & $2.60\times 10^{-15}$ & 23 \\
        $45°$ & $1.59\times 10^{0}$   & $7.30\times 10^{-1}$  & 33 \\
        $60°$ & $8.48\times 10^{-1}$  & $7.02\times 10^{-1}$  & 33 \\
        $90°$ & $2.57\times 10^{-15}$ & $2.74\times 10^{-15}$ & 38 \\
        \bottomrule
    \end{tabular}
\end{table}

从表~\ref{tab:angle}可以看出，当旋转角度不超过$30°$时，ICP算法能够精确恢复真实变换，旋转误差和平移误差均在$10^{-15}$量级；当旋转角度为$45°$和$60°$时，出现了较大的配准误差，说明ICP在这些角度下陷入了局部最优；而旋转角度为$90°$时又恢复了精确配准。这一现象可以从点云的几何结构来理解：模拟兔子形状的点云具有非对称的几何特征，在某些旋转角度下，最近点对应关系可能偏离真实对应，导致ICP收敛到错误的局部最小值。小角度（$\leq 30°$）下的平均迭代次数为19.7次，大角度（$\geq 30°$）下的平均迭代次数为34.7次，表明大角度旋转需要更多的迭代才能收敛。这一结果揭示了ICP算法的一个重要局限性：它仅保证收敛到局部最小值，对于大旋转角度的情况，可能需要结合粗配准方法来提供良好的初始位姿。

\section{实验四：不同重叠率下的ICP配准}

在实际应用中，两组点云往往只有部分重叠，因此研究重叠率对ICP配准精度的影响具有重要意义。本实验通过随机采样源点云和目标点云的不同子集来模拟不同的重叠率，重叠率从100\%递减至50\%。表~\ref{tab:overlap}列出了不同重叠率下的配准结果。

\begin{table}[H]
    \centering
    \caption{不同重叠率下的ICP配准结果}
    \label{tab:overlap}
    \begin{tabular}{ccc}
        \toprule
        重叠率 & 旋转误差 & 平移误差 \\
        \midrule
        100\% & $2.61\times 10^{-15}$ & $2.50\times 10^{-15}$ \\
        90\%  & $3.69\times 10^{-3}$  & $2.21\times 10^{-3}$  \\
        80\%  & $1.95\times 10^{-2}$  & $7.22\times 10^{-3}$  \\
        70\%  & $1.43\times 10^{0}$   & $6.84\times 10^{-1}$  \\
        60\%  & $1.90\times 10^{-2}$  & $3.00\times 10^{-2}$  \\
        50\%  & $5.99\times 10^{-2}$  & $1.16\times 10^{-2}$  \\
        \bottomrule
    \end{tabular}
\end{table}

从表~\ref{tab:overlap}可以看出，当重叠率为100\%时，配准误差在机器精度范围内；当重叠率降至90\%和80\%时，配准误差虽有所增大但仍保持在较小水平；当重叠率为70\%时，出现了较大的配准误差，说明此时ICP陷入了局部最优；而在60\%和50\%的重叠率下，误差又有所回落，但相比完全重叠的情况仍有明显增大。总体而言，重叠率的降低会导致最近点对应关系中错误匹配的比例增加，从而影响SVD求解的刚体变换的准确性。ICP算法对部分重叠具有一定的鲁棒性，但当重叠率过低时，配准精度会显著下降。

\section{SVD算法正确性验证}

除了ICP配准实验外，本文还对手动实现的SVD分解进行了独立的正确性验证。测试了四种不同类型的矩阵：$3\times 3$方阵、$4\times 3$矩形矩阵、$3\times 5$矩形矩阵和秩亏矩阵。验证指标包括重构误差$\|\bm{A}-\bm{U}\bm{\Sigma}\bm{V}^{\mathsf{T}}\|$、$\bm{U}$的正交性误差$\|\bm{U}^{\mathsf{T}}\bm{U}-\bm{I}\|$、$\bm{V}$的正交性误差$\|\bm{V}^{\mathsf{T}}\bm{V}-\bm{I}\|$以及奇异值与NumPy库的对比误差。表~\ref{tab:svd}列出了验证结果。

\begin{table}[H]
    \centering
    \caption{SVD分解正确性验证结果}
    \label{tab:svd}
    \begin{tabular}{ccccc}
        \toprule
        测试矩阵 & 重构误差 & $\bm{U}$正交性误差 & $\bm{V}$正交性误差 & 奇异值对比误差 \\
        \midrule
        方阵$3\times 3$     & $1.75\times 10^{-15}$ & $1.08\times 10^{-15}$ & $1.28\times 10^{-15}$ & $3.51\times 10^{-16}$ \\
        矩形$4\times 3$     & $7.97\times 10^{-16}$ & $2.29\times 10^{-15}$ & $3.50\times 10^{-16}$ & $9.04\times 10^{-16}$ \\
        矩形$3\times 5$     & $5.56\times 10^{-16}$ & $4.04\times 10^{-16}$ & $7.16\times 10^{-16}$ & $8.38\times 10^{-16}$ \\
        秩亏矩阵$3\times 3$ & $4.56\times 10^{-15}$ & $1.64\times 10^{-16}$ & $4.68\times 10^{-16}$ & $3.78\times 10^{-15}$ \\
        \bottomrule
    \end{tabular}
\end{table}

从表~\ref{tab:svd}可以看出，所有测试用例的重构误差均在$10^{-15}$量级，正交性误差也在$10^{-15}$量级，奇异值与NumPy库的对比误差在$10^{-16}$至$10^{-15}$量级，均在机器精度范围内。这充分验证了手动实现的SVD分解算法的正确性，包括基于$\bm{A}^{\mathsf{T}}\bm{A}$或$\bm{A}\bm{A}^{\mathsf{T}}$特征分解构造奇异向量的方法以及利用改进Gram-Schmidt方法补全正交基的策略。

\begin{figure}[H]
    \centering
    \includegraphics[width=0.95\textwidth]{svd_verification.png}
    \caption{SVD正确性验证：奇异值对比与不同尺寸矩阵的重构误差}
    \label{fig:svd_verification}
\end{figure}

图~\ref{fig:svd_verification}左图展示了手动实现与NumPy库计算的奇异值对比，两者完全一致；右图展示了不同尺寸矩阵的SVD重构误差，均在$10^{-15}$量级，进一步证实了SVD实现的正确性。

\begin{figure}[H]
    \centering
    \includegraphics[width=0.95\textwidth]{icp_convergence.png}
    \caption{ICP收敛曲线：无噪声、不同噪声水平和不同旋转角度}
    \label{fig:convergence}
\end{figure}

图~\ref{fig:convergence}展示了ICP算法在不同条件下的收敛曲线。左图为无噪声条件下的收敛过程，均方误差从初始值快速下降至接近零；中图为不同噪声水平下的收敛曲线，噪声越大，最终收敛的误差水平越高，但收敛速度差异不大；右图为不同旋转角度下的收敛曲线，小角度旋转收敛较快且最终精度高，大角度旋转收敛较慢且可能陷入局部最优。

\begin{figure}[H]
    \centering
    \includegraphics[width=0.95\textwidth]{icp_error_analysis.png}
    \caption{噪声、旋转角度和重叠率对ICP配准精度的影响}
    \label{fig:error_analysis}
\end{figure}

图~\ref{fig:error_analysis}综合展示了噪声、旋转角度和重叠率三个因素对ICP配准精度的影响。左图显示旋转误差和平移误差均随噪声水平近似线性增长；中图显示小角度旋转的配准精度远高于大角度旋转；右图显示重叠率降低导致配准精度下降，且在某些重叠率下可能出现局部最优问题。

% ============================================================
% 第五章 结果分析与讨论
% ============================================================
\chapter{结果分析与讨论}

\section{模型正确性验证}

通过四组仿真实验，本文从多个维度验证了所构建的SVD和ICP模型的正确性。首先，SVD分解的正确性通过独立测试得到了充分验证：四种不同类型的矩阵（方阵、两种矩形矩阵和秩亏矩阵）的重构误差均在$10^{-15}$量级，与NumPy库的计算结果完全一致，这证明了基于$\bm{A}^{\mathsf{T}}\bm{A}$特征分解构造右奇异向量、利用$\bm{u}_i=\bm{A}\bm{v}_i/\sigma_i$构造左奇异向量以及改进Gram-Schmidt补全正交基的方法论是正确的。其次，ICP算法的正确性通过无噪声配准实验得到了验证：在没有任何扰动的情况下，ICP算法精确恢复了预设的刚体变换参数，旋转矩阵误差为$2.77\times 10^{-15}$，平移向量误差为$5.98\times 10^{-16}$，估计旋转矩阵的正交性和行列式约束均得到满足，这证明了基于SVD的刚体变换求解方法和ICP迭代框架的实现是正确的。

\section{噪声鲁棒性分析}

噪声鲁棒性实验表明，ICP算法的配准误差随噪声水平近似线性增长，旋转误差的增长斜率约为0.543，平移误差的增长斜率约为0.564。这一线性关系可以从理论上加以解释：在ICP的每次迭代中，噪声导致最近点搜索产生偏差，使得部分点对对应关系偏离真实对应；SVD虽然给出了当前对应关系下的最优刚体变换，但由于对应关系本身存在误差，最终求得的变换参数也会偏离真实值。噪声水平越高，对应关系的偏差越大，配准误差也就越大。在$\sigma=0.2$的较高噪声水平下，旋转误差为0.106，平移误差为0.108，相对于点云的尺度（约2个单位）而言，这一误差是可以接受的，说明ICP算法对中等程度的噪声具有良好的鲁棒性。

\section{旋转角度敏感性分析}

旋转角度实验揭示了ICP算法的一个重要特性：对于具有丰富几何特征的非对称点云，ICP在小角度旋转（$\leq 30°$）下能够精确收敛到全局最优解，但在某些大角度（如$45°$和$60°$）下可能陷入局部最优。这一现象是ICP算法的固有局限性，而非实现错误。ICP算法的收敛性保证仅针对局部最小值：在每次迭代中，最近点搜索和SVD求解都使目标函数不增，但目标函数可能存在多个局部最小值，初始位姿距离全局最优越远，陷入局部最优的风险就越大。值得注意的是，$90°$旋转反而能够精确配准，这是因为在该特定角度下，点云的最近点对应关系恰好与真实对应一致，使得ICP能够跳出局部最小值。这一结果提示我们，在实际应用中，当初始旋转角度较大时，应当结合粗配准方法（如基于特征描述子的匹配或随机采样一致性）来提供良好的初始位姿，以提高ICP配准的可靠性。

\section{重叠率依赖性分析}

重叠率实验表明，ICP算法的配准精度对点云重叠率有较强的依赖性。当重叠率为100\%时，配准误差在机器精度范围内；当重叠率降至90\%和80\%时，误差虽有所增大但仍保持在较小水平；当重叠率进一步降低时，配准精度出现明显波动，在某些重叠率下甚至陷入局部最优。这是因为重叠率降低意味着源点云中有更多点无法在目标点云中找到正确的对应点，这些错误对应会干扰SVD求解的刚体变换。此外，由于本实验中源点云和目标点云的子集是独立随机采样的，不同重叠率下的点云组成存在差异，这也导致了配准精度的波动。总体而言，ICP算法对部分重叠具有一定的鲁棒性，但在重叠率低于70\%时，配准结果的可靠性会显著下降。

\section{SVD在ICP中的核心作用}

综合以上分析，SVD在ICP算法中扮演着不可或缺的核心角色。在ICP的每次迭代中，基于当前最近点对应关系求解最优刚体变换是算法的关键步骤，而SVD提供了一种高效、稳定且解析的求解方法。通过将互协方差矩阵$\bm{H}=\bm{P}'^{\mathsf{T}}\bm{Q}'$进行SVD分解$\bm{H}=\bm{U}\bm{\Sigma}\bm{V}^{\mathsf{T}}$，最优旋转矩阵可以直接写为$\bm{R}=\bm{V}\bm{U}^{\mathsf{T}}$，最优平移向量为$\bm{t}=\bar{\bm{q}}-\bm{R}\bar{\bm{p}}$，无需进行迭代优化或梯度计算。这一解析解不仅计算效率高，而且数值稳定性好，因为SVD分解本身具有优良的数值性质。正是由于SVD的这一特性，ICP算法才能在每次迭代中快速且准确地求解最优变换，从而保证了算法的收敛性和实用性。

% ============================================================
% 第六章 结论
% ============================================================
\chapter{结论}

本文对奇异值分解和迭代最近点算法进行了系统的数学推导、编程实现和仿真实验验证，主要工作和结论如下。

在算法推导方面，本文从矩阵特征分解的角度严格推导了SVD的存在性和构造方法，详细阐述了基于$\bm{A}^{\mathsf{T}}\bm{A}$或$\bm{A}\bm{A}^{\mathsf{T}}$特征分解构造奇异向量的理论依据，以及利用改进Gram-Schmidt方法补全正交基的策略。在此基础上，推导了ICP算法中基于SVD求解最优刚体变换的完整过程，包括平移向量的解耦、互协方差矩阵的构造、最优旋转矩阵的解析求解以及反射修正，并给出了ICP的完整迭代流程和收敛性分析。

在仿真实验方面，本文设计了四组实验来验证算法的正确性和分析算法的性能。无噪声配准实验表明，ICP算法能够精确恢复真实的刚体变换参数，旋转矩阵误差和平移向量误差均在$10^{-15}$量级，验证了算法实现的正确性。噪声鲁棒性实验表明，配准误差随噪声水平近似线性增长，ICP对中等程度的噪声具有良好的鲁棒性。旋转角度敏感性实验表明，小角度旋转下ICP能够精确收敛，大角度旋转可能陷入局部最优，这是ICP算法的固有局限性。重叠率依赖性实验表明，ICP对部分重叠具有一定鲁棒性，但重叠率过低会导致配准精度下降。

在模型正确性验证方面，手动实现的SVD分解在四种不同类型矩阵上的重构误差均在机器精度范围内，与NumPy库的计算结果完全一致；ICP算法在无噪声条件下精确恢复了真实变换参数，估计旋转矩阵满足正交性和行列式约束。这些结果充分验证了所构建的SVD和ICP模型的正确性。
\end{document}
